\documentclass[11pt,a4paper]{report}

\usepackage{fontspec}
\usepackage{amsmath,amssymb}
\usepackage{listings}
\usepackage{xcolor}
\usepackage{hyperref}
\usepackage{geometry}
\usepackage{longtable}
\usepackage{booktabs}
\usepackage{enumitem}
\usepackage{parskip}
\usepackage{titlesec}

\geometry{margin=2.7cm}

\definecolor{codebg}{RGB}{245,245,245}
\definecolor{codegray}{RGB}{110,110,110}
\definecolor{codeblue}{RGB}{30,80,160}
\definecolor{codegreen}{RGB}{30,120,60}

\lstdefinestyle{py}{
  backgroundcolor=\color{codebg},
  basicstyle=\ttfamily\small,
  keywordstyle=\color{codeblue}\bfseries,
  commentstyle=\color{codegray}\itshape,
  stringstyle=\color{codegreen},
  numbers=none,
  breaklines=true,
  showstringspaces=false,
  frame=single,
  framerule=0.3pt,
  rulecolor=\color{codegray},
  xleftmargin=1em,
  xrightmargin=1em,
  aboveskip=8pt,
  belowskip=8pt,
  language=Python,
}
\lstset{style=py}

\lstdefinestyle{yaml}{
  style=py,
  language=,
  keywordstyle=\color{codeblue}\bfseries,
}

\hypersetup{
  colorlinks=true,
  linkcolor=codeblue,
  urlcolor=codeblue,
  citecolor=codeblue,
  bookmarksopen=true,
}

\titleformat{\chapter}[hang]{\normalfont\huge\bfseries}{\thechapter.}{20pt}{\huge}
\setlength{\parindent}{0pt}
\setlength{\parskip}{6pt}

\newcommand{\arq}[1]{\texttt{#1}}
\newcommand{\code}[1]{\texttt{#1}}

\title{\Huge Documentação Técnica do Projeto FEMa\\[6pt]
\Large Comparação de Bases de Interpolação para o Finite Element Machine}
\author{}
\date{\today}

\begin{document}

\maketitle
\tableofcontents

\chapter{Introdução}

\section{Objetivo científico do projeto}

Este projeto implementa e avalia o \textbf{FEMa} (\emph{Finite Element Machine}), um método de
classificação/regressão baseado em interpolação por vizinhos próximos: para prever o rótulo (ou
valor) de uma amostra nova, o FEMa busca os $k$ vizinhos mais próximos no conjunto de treino e
combina os valores desses vizinhos usando pesos derivados de uma \textbf{função de base}
$\varphi(d)$, onde $d$ é a distância até cada vizinho.

O objeto de estudo central do projeto \textbf{não é o FEMa em si}, mas sim
\textbf{a comparação entre as diferentes funções de base} que podem ser usadas dentro dele —
Shepard, RBF Gaussiana, Multiquadrática, Wendland C2, splines, kernels de Cauchy, Student-t,
Racional Quadrático, entre outras (21 ao todo, listadas na Seção~\ref{sec:factory-basis}). A
pergunta de pesquisa é, essencialmente: \emph{dado um dataset e uma tarefa (classificação ou
regressão), qual função de base produz o melhor desempenho preditivo, e por quê?}

Essa é a razão pela qual, como será detalhado ao longo deste documento, toda a arquitetura do
código — do desenho dos \emph{plugins} de modelo à organização dos resultados em disco — foi
construída para que a \textbf{base de interpolação seja o eixo estrutural de comparação}, e não
apenas mais um hiperparâmetro entre outros. Baselines externos (regressão logística, k-NN) existem
apenas como \emph{referência metodológica} — para responder ``o FEMa com a melhor base é
competitivo com um classificador simples?'' — e ficam deliberadamente isolados do restante da
comparação.

\section{Visão geral da arquitetura}

O projeto segue uma arquitetura em camadas, na qual cada diretório do repositório tem uma
responsabilidade única e bem delimitada. A regra geral é: \textbf{uma camada só conhece a camada
imediatamente abaixo dela}. Isso significa, por exemplo, que a camada de persistência
(\arq{persistence/}) não sabe o que é uma ``base'' ou um ``dataset'' — ela só grava e lê JSON em
caminhos genéricos; quem decide o \emph{significado} de cada segmento do caminho é a camada de
orquestração (\arq{pipeline/}).

O fluxo de dependências entre camadas, da base para o topo, resumido a seguir:

\begin{enumerate}[label=\arabic*.]
  \item \textbf{\arq{core/}} — o núcleo matemático do FEMa (bases, distâncias, busca de vizinhos,
        os modelos \code{FEMaClassifier}/\code{FEMaRegressor}). É tratado como uma
        \emph{biblioteca fechada}: nenhuma decisão de experimento deveria exigir alterá-lo.
  \item \textbf{\arq{datasets/}} — carregamento e split de dados, desacoplado de qualquer modelo.
  \item \textbf{\arq{models/}} — a camada de \emph{plugins}: adapta o FEMa (\arq{core/}) e os
        baselines externos (scikit-learn) a um contrato único (\code{ModelPlugin}).
  \item \textbf{\arq{metrics/}}, \textbf{\arq{tuning/}}, \textbf{\arq{training/}},
        \textbf{\arq{evaluation/}} — blocos de construção independentes usados pelo orquestrador:
        cálculo de métricas, geração de combinações de hiperparâmetros, treino e avaliação.
  \item \textbf{\arq{persistence/}} — escrita atômica dos resultados em JSON e sistema de
        checkpoint (retomada de execuções interrompidas).
  \item \textbf{\arq{pipeline/run\_model.py}} — o orquestrador único que liga todas as camadas
        acima em um fluxo completo: carregar dados $\to$ checkpoint $\to$ busca de
        hiperparâmetros $\to$ validação cruzada $\to$ resumo $\to$ retreino final $\to$ avaliação
        no teste.
  \item \textbf{\arq{reporting/}} — lê os resultados já persistidos e produz tabelas, JSON
        consolidado e gráficos comparando as bases entre si.
  \item \textbf{\arq{utils/}} — utilitários transversais (log, seed, tempo, hash).
  \item \textbf{\arq{configs/}} — arquivos YAML declarativos que parametrizam experimentos, sem
        exigir alteração de código.
  \item \textbf{\arq{main.py}} — ponto de entrada de linha de comando.
\end{enumerate}

Os capítulos a seguir documentam cada camada nessa mesma ordem — da matemática para a
orquestração —, e o Capítulo~\ref{ch:expansao} explica concretamente como estender o projeto em
cada uma das direções mais prováveis: novos datasets, novas bases, análise de regressão e
comparação com outros modelos.

\chapter{O núcleo matemático: \arq{core/}}
\label{ch:core}

\arq{core/} contém tudo o que é \textbf{matemática do FEMa em si} — as funções de base, as
métricas de distância, a estrutura de busca de vizinhos e os dois modelos
(\code{FEMaClassifier}, \code{FEMaRegressor}). Por convenção do projeto, esta pasta é tratada como
uma biblioteca estável: decisões de \emph{experimento} (que faixa de hiperparâmetros testar, quais
métricas calcular, como organizar os resultados) vivem fora dela, nas camadas superiores.

\section{\arq{core/math/distances/} — métricas de distância}

Três arquivos pequenos definem o contrato de distância:

\begin{itemize}
  \item \arq{base\_distance.py}: classe abstrata \code{BaseDistance} com um único método,
        \code{compute(x1, x2)}, vetorizado (aceita tanto um par de vetores quanto um vetor contra
        uma matriz de N vetores).
  \item \arq{euclidean\_distance.py}: \code{EuclideanDistance}, a métrica usada em praticamente
        todo o projeto — norma $L_2$ da diferença, com \code{np.linalg.norm}.
  \item \arq{manhattan\_distance.py}: \code{ManhattanDistance}, norma $L_1$, disponível mas não
        usada em nenhum experimento atual (ver Seção~\ref{sec:expansao-modelos} sobre como
        conectá-la).
\end{itemize}

\section{\arq{core/math/neighboor\_search/} — busca de vizinhos}

\code{BaseSearch} define o contrato: \code{build(X)} indexa o conjunto de treino, e
\code{query(sample, k)} retorna os índices e distâncias dos $k$ vizinhos mais próximos de uma
amostra (ou de todos, se $k=0$).

A única implementação existente, \code{BruteForceSearch} (\arq{brute\_force.py}), calcula a
distância de \code{sample} contra \textbf{todos} os pontos de treino e usa
\code{np.argpartition} para selecionar os $k$ menores em tempo $O(n)$ (sem ordenar o vetor
inteiro). É deliberadamente simples: para os tamanhos de dataset usados no projeto, uma estrutura
de indexação espacial (KD-tree, ball-tree) traria complexidade adicional sem necessidade prática.
Trocar a estratégia de busca por uma mais eficiente é uma extensão natural e isolada (basta
implementar \code{BaseSearch} — ver Seção~\ref{sec:expansao-modelos}).

\section{\arq{core/math/basis/} — as funções de base}
\label{sec:factory-basis}

Esta é a parte mais extensa e mais importante de \arq{core/}, pois \textbf{as funções de base são
o objeto de estudo do projeto}.

\subsection{\arq{parameters.py} — o contêiner único de hiperparâmetros}

Antes de qualquer base específica, vale entender \code{BasisParameters}
(\arq{parameters.py}): um \emph{dataclass} único com nove campos opcionais
(\code{z, epsilon, c, alpha, beta, nu, h, l, p}), todos com valor-padrão \code{None}. Essa é uma
decisão de design central: em vez de cada base ter sua própria assinatura de função (o que
obrigaria quem chama a base — o classificador, o harness de validação, a camada de tuning — a
conhecer a fórmula de cada uma para montar os argumentos certos), \textbf{toda base recebe a mesma
assinatura} \code{compute\_weights(dists, params)}, e cada implementação lê, de \code{params},
apenas os campos que sua fórmula usa. Um campo faz o papel de vários símbolos matemáticos
diferentes dependendo da base (por exemplo, \code{c} é usado por Multiquadrática, Inversa
Multiquadrática, Logarítmica e Sigmoidal, cada uma com um papel distinto na fórmula) — isso é
intencional: o nome é só o ``slot'' de configuração, não implica semântica compartilhada.

\subsection{\arq{base\_basis.py} — a classe abstrata \code{BaseBasis}}

\code{BaseBasis} separa duas camadas deliberadamente:

\begin{itemize}
  \item \code{evaluate(dists, params)}: o valor \textbf{cru} (não normalizado) de $\varphi(d)$.
        Cada subclasse concreta implementa apenas isto — a fórmula da base em si.
  \item \code{compute\_weights(dists, params)}: o peso \textbf{normalizado} (partição da
        unidade), $w = \varphi(d) / \sum \varphi(d)$. Implementado uma única vez, na classe base,
        em cima de \code{evaluate()}. Nenhuma subclasse deveria reimplementar isso.
\end{itemize}

Essa separação corrige, por construção, uma classe inteira de bug que existia em versões
anteriores do código: a normalização feita incorretamente dentro de cada base individualmente
(nos comentários do projeto, ``\code{weights /= weights}'' em vez de
``\code{weights /= weights.sum()}''). Centralizar a normalização em \code{BaseBasis} elimina essa
possibilidade de erro para qualquer base nova.

Além disso, se $\sum \varphi(d) = 0$ (uma base de suporte compacto em que nenhum vizinho cai
dentro do raio, por exemplo), \code{compute\_weights} recorre a pesos uniformes em vez de propagar
um \code{NaN} silenciosamente.

Cada subclasse declara dois atributos de classe:

\begin{itemize}
  \item \code{PARAMS}: os nomes dos campos de \code{BasisParameters} que a fórmula
        \textbf{exige}. O método \code{\_require(params)} extrai esses campos e levanta um
        \code{ValueError} claro se algum estiver \code{None}, em vez de deixar o erro estourar
        mais adiante como uma exceção aritmética difícil de rastrear.
  \item \code{OPTIONAL\_PARAMS}: campos que a base usa, mas que são opcionais (o caso mais comum
        é \code{h} nas bases de suporte compacto, onde \code{None} significa ``raio automático,
        igual à maior distância observada'').
\end{itemize}

\subsection{\arq{\_compactSupport.py} — bases de suporte compacto}

\code{CompactSupportBasis} é uma classe intermediária, herdada por Wendland C2, Cubic Spline,
Quartic Spline e Cosine. Fornece \code{\_normalized\_distance(dists, h)}, que calcula
$r = d/h$ — usando $h = \max(d)$ automaticamente quando \code{h} não é informado. Fora do
intervalo $r \le 1$, a base vale zero: são funções de \textbf{suporte compacto}, que ignoram
vizinhos além de um certo raio.

\subsection{As 21 bases registradas}

A Tabela~\ref{tab:bases} resume as 21 funções de base atualmente implementadas, todas em
\arq{core/math/basis/}, registradas na fábrica \code{Basis} (\arq{factory\_basis.py},
Seção~\ref{sec:factory}). Note que várias bases têm notas de equivalência matemática
documentadas no próprio código-fonte (docstrings marcadas ``NOTA''): decisões de nomenclatura e
de fórmula ainda pendentes de confirmação com o orientador, deixadas explícitas em vez de
escondidas.

\begingroup
\small
\begin{longtable}{@{}p{2.6cm}p{6.2cm}p{4.6cm}@{}}
\caption{As 21 funções de base registradas em \code{Basis} (\arq{factory\_basis.py}).}
\label{tab:bases} \\
\toprule
\textbf{Nome (registry)} & \textbf{Fórmula} $\varphi(d)$ & \textbf{Parâmetros} \\
\midrule
\endfirsthead
\toprule
\textbf{Nome (registry)} & \textbf{Fórmula} $\varphi(d)$ & \textbf{Parâmetros} \\
\midrule
\endhead
\code{shepard} & $1/d^z$ & \code{z} \\
\code{radial} & $\exp(-0.5\,(d/r_0)^2)$, $r_0=z$ & \code{z} \\
\code{rbf\_gaussian} & $\exp(-(\epsilon d)^2)$ & \code{epsilon} \\
\code{multiquadratic} & $\sqrt{d^2+c^2}$ (crescente em $d$) & \code{c} \\
\code{inverse\_multiquadratic} & $1/\sqrt{d^2+c^2}$ & \code{c} \\
\code{wendland\_c2} & $(1-r)^4(4r+1)$, $r=d/h$, $0$ se $r>1$ & \code{h} (opc.) \\
\code{cubic\_spline} & $(1-r)^4$, $r=d/h$, $0$ se $r>1$ & \code{h} (opc.) \\
\code{quartic\_spline} & $(1-r)^4$, $r=d/h$, $0$ se $r>1$ (idêntica a \code{cubic\_spline} — nota de decisão pendente) & \code{h} (opc.) \\
\code{gen\_exponential} & $\exp(-\epsilon\,d^p)$ & \code{epsilon}, \code{p} \\
\code{softmax\_radial} & $\exp(-\beta d)$ (nota: equivalente a \code{laplacian} após normalização) & \code{beta} \\
\code{attention} & $1/(1+d^2)$ (nota: equivalente a \code{cauchy} com $\epsilon=1$) & --- \\
\code{logarithmic} & $1/\log(1+d+c)$ & \code{c} \\
\code{harmonic} & $(1+d^2/\nu)^{-(\nu+1)/2}$ (nota: formalmente idêntica a \code{student\_t}) & \code{nu} \\
\code{laplacian} & $\exp(-\epsilon d)$ & \code{epsilon} \\
\code{cauchy} & $1/(1+(\epsilon d)^2)$ & \code{epsilon} \\
\code{student\_t} & $(1+d^2/\nu)^{-(\nu+1)/2}$ (kernel-t, van der Maaten \& Hinton, 2008) & \code{nu} \\
\code{cosine} & $\cos^2(\pi r/2)$, $r=d/h$, $0$ se $r>1$ & \code{h} (opc.) \\
\code{sigmoidal} & $1/(1+\exp(\alpha(d-c)))$ & \code{alpha}, \code{c} \\
\code{lorentzian} & $1/(1+d^4)$ (nota: na literatura ``Lorentzian'' costuma ser sinônimo de Cauchy) & --- \\
\code{entropic} & $\exp(-\beta d^2)$ & \code{beta} \\
\code{rational\_quadratic} & $(1+d^2/(2\alpha l^2))^{-\alpha}$ (Rasmussen \& Williams, 2006) & \code{alpha}, \code{l} \\
\bottomrule
\end{longtable}
\endgroup

Todas seguem o mesmo padrão de arquivo — pequeno, uma classe, um método \code{evaluate}:

\begin{lstlisting}
class ShepardBasis(BaseBasis):
    PARAMS = ("z",)

    def evaluate(self, dists, params):
        z = self._require(params)["z"]
        dists = np.where(dists == 0, DELTA, dists)  # evita divisao por zero
        return 1.0 / (dists ** z)
\end{lstlisting}

\subsection{\arq{factory\_basis.py} — a fábrica \code{Basis}}
\label{sec:factory}

\code{Basis} é o ponto único de acesso às implementações: \code{Basis.get(nome, search)}
instancia a base pelo nome (string), e \code{Basis.available()} retorna a lista ordenada de todos
os nomes registrados no dicionário interno \code{\_CLASSES}. Essa função é usada extensivamente
por outras camadas (o harness de validação matemática, os \emph{plugins} de modelo, o módulo de
comparação de bases) para \textbf{descobrir automaticamente} quais bases existem, sem duplicar
essa lista em vários lugares do projeto — uma base nova só precisa de uma entrada em
\code{\_CLASSES} para se tornar visível a todo o resto do sistema (ver
Seção~\ref{sec:expansao-base}).

\section{\arq{core/models/} — os modelos FEMa}

\subsection{\arq{base\_model.py} — \code{FEMaBaseModel}}

Classe abstrata que define o contrato comum: o modelo guarda \code{X\_train}/\code{y\_train},
delega a indexação ao \code{search} injetado no construtor, e implementa \code{fit}/\code{predict}
como métodos abstratos. A ideia de responsabilidade é: \emph{o modelo orquestra search + basis e
faz o produto interno final; a busca de vizinhos e o cálculo de pesos são delegados}.

\subsection{\arq{fema\_classifier.py} — \code{FEMaClassifier}}

Para classificação, o \code{fit} constrói uma matriz de probabilidades por classe: uma linha por
classe, com $1.0$ nas posições de treino cuja classe é aquela e $0.0$ nas demais (codificação
one-hot transposta). O \code{predict}, para cada amostra de teste, busca os $k$ vizinhos uma
única vez, calcula os pesos com a base escolhida, e então interpola a probabilidade de
\textbf{cada} classe usando os \textbf{mesmos} pesos (produto interno entre o vetor de pesos e a
linha da classe). O rótulo previsto é o \code{argmax} das probabilidades interpoladas.

\subsection{\arq{fema\_regressor.py} — \code{FEMaRegressor}}

Para regressão, o processo é mais direto: o \code{fit} apenas guarda \code{X\_train}/\code{y\_train}
e indexa \code{X\_train} no \code{search}; o \code{predict} interpola diretamente o valor-alvo
(produto interno entre pesos e \code{y\_train[indices]}), com uma salvaguarda para o caso raro em
que a predição resulte em \code{NaN} (cai para a média de \code{y\_train}).

Este componente já existe e está pronto para uso — ver Seção~\ref{sec:expansao-regressao} sobre
como habilitar experimentos de regressão de ponta a ponta.

\chapter{Camada de dados: \arq{datasets/}}

\section{\arq{registry.py} — o que existe e onde}

Define \code{DatasetSpec}, um \emph{dataclass} imutável com o caminho do CSV, a coluna-alvo, as
colunas a descartar e o delimitador. O dicionário \code{\_REGISTRY} mapeia um nome lógico de
dataset para essa especificação. Quatro datasets estão registrados hoje:

\begin{itemize}
  \item \code{fetal\_health} — CSV real, alvo \code{fetal\_health} (3 classes).
  \item \code{iris} — CSV real, alvo \code{Species}, descarta a coluna \code{Id}.
  \item \code{classification\_data} — CSV real com delimitador \code{;}, alvo \code{class}.
  \item \code{synthetic\_demo} — não tem arquivo; é gerado em tempo de execução (ver
        \arq{loader.py}) — usado nos testes automatizados, por ser rápido e não depender de
        arquivos externos.
\end{itemize}

Este módulo \textbf{nunca lê nem transforma dados} — apenas descreve onde/como encontrá-los. A
função \code{get\_dataset\_spec(name)} valida que o arquivo existe no disco antes de devolver a
especificação (levanta \code{FileNotFoundError} cedo, com mensagem clara, em vez de falhar mais
adiante dentro do \emph{pandas}).

\section{\arq{loader.py} — leitura, split e padronização}

\code{load\_dataset(dataset\_name, seed, val\_ratio, test\_ratio, scale)} é a única função pública
relevante. Faz, nesta ordem:

\begin{enumerate}
  \item Lê o CSV (via \code{pandas.read\_csv}, normalizando nomes de coluna e removendo o BOM
        \code{\textbackslash ufeff} que alguns arquivos trazem) ou gera o dataset sintético
        (\code{sklearn.datasets.make\_classification}).
  \item Separa \code{X} (todas as colunas menos o alvo, convertidas para \code{float64}) de
        \code{y} (convertido para inteiro diretamente se já for numérico, ou por
        \code{codes} de categoria caso contrário).
  \item Faz o split treino/validação/teste em duas chamadas encadeadas de
        \code{train\_test\_split} (ambas \code{stratify}), de forma \textbf{determinística}
        (mesma \code{seed} sempre produz o mesmo split).
  \item Se \code{scale=True} (padrão), padroniza todas as partições
        ($x' = (x-\mu)/\sigma$) usando \textbf{apenas} estatísticas do conjunto de treino — esta
        padronização é essencial para o FEMa, que é um método baseado em distância: sem ela,
        variáveis em escalas diferentes dominariam artificialmente o cálculo de distância.
\end{enumerate}

O retorno é um \code{LoadedData} (\emph{dataclass} imutável) com
\code{X\_train/y\_train/X\_val/y\_val/X\_test/y\_test} e o nome do dataset. Esta camada não conhece
nenhum modelo nem faz parte do pipeline de tuning/avaliação — é puramente sobre dados.

\chapter{Camada de modelos: \arq{models/}}
\label{ch:models}

\section{\arq{base.py} — o contrato \code{ModelPlugin}}

Esta é a peça de desenho mais importante do projeto depois de \arq{core/}: \textbf{nenhuma outra
camada importa um algoritmo diretamente}. Tudo — treino, avaliação, tuning, persistência — fala
apenas com a interface abstrata \code{ModelPlugin}. Isso é o que permite trocar de modelo (ou
adicionar um novo) sem tocar em \arq{pipeline/}, \arq{evaluation/}, \arq{training/} ou
\arq{tuning/}.

Os métodos do contrato:

\begin{itemize}
  \item \code{create\_model(params, random\_state)} \emph{(abstrato)}: instancia (sem treinar) um
        estimador com os hiperparâmetros dados.
  \item \code{parameter\_grid()} \emph{(abstrato)}: espaço de busca \textbf{discreto}, para
        \code{grid\_search} — formato \code{\{nome: [valores]\}}.
  \item \code{random\_search\_space()} \emph{(opcional, com \code{default} = reaproveitar
        \code{parameter\_grid()})}: espaço de busca \textbf{contínuo}, para \code{random\_search}
        — aceita também tuplas \code{("uniform", low, high)}, \code{("loguniform", low, high)} e
        \code{("randint", low, high)}. Esta é uma extensão em relação ao projeto de referência que
        inspirou esta arquitetura (que só tinha grades discretas): necessária porque o FEMa e os
        baselines aqui têm hiperparâmetros contínuos (\code{z}, \code{epsilon}, \code{C}, \ldots)
        mal representados por uma grade fixa.
  \item \code{fit(estimator, X, y, X\_val, y\_val)}: por padrão, chama \code{estimator.fit(X, y)};
        plugins podem sobrescrever se precisarem de \code{X\_val}/\code{y\_val} (ex.: \emph{early
        stopping}).
  \item \code{predict}/\code{predict\_proba}: delegam ao estimador; \code{predict\_proba} retorna
        \code{None} graciosamente se o modelo não suportar (o cálculo de métricas trata isso sem
        quebrar, pulando métricas que dependem de probabilidade, como \code{roc\_auc}).
  \item \code{save}/\code{load}: serialização via \code{joblib}.
\end{itemize}

\section{\arq{fema.py} — o FEMa como (contexto, base)}
\label{sec:models-fema}

Este é o arquivo que mais evoluiu ao longo do desenvolvimento do projeto, e a evolução em si é
instrutiva. Numa versão anterior, existia um único \code{FEMaPlugin} ``genérico'', e a base de
interpolação era apenas mais um valor dentro do espaço de hiperparâmetros buscado
(\code{basis\_function}), junto com \code{k}, \code{z}, \code{epsilon} etc. Isso tinha um problema
conceitual sério: o \emph{ranking} final do resumo (\arq{persistence/summary\_builder.py}) acabava
misturando combinações de bases \textbf{diferentes} na mesma lista, de modo que a ``melhor
combinação'' encontrada respondia à pergunta errada (``qual é a melhor combinação global entre
todas as bases misturadas'', em vez de ``qual é o melhor desempenho \emph{de cada} base'').

A arquitetura atual resolve isso fixando o \textbf{contexto} (\code{"classifier"} ou
\code{"regressor"}) e a \textbf{base} na própria instância do plugin:

\begin{lstlisting}
class FEMaPlugin(ModelPlugin):
    def __init__(self, context: str, basis: str):
        ...
        self.context = context
        self.basis = basis
        self.name = f"fema:{context}:{basis}"
        self.supports_proba = context == "classifier"

    def parameter_grid(self) -> Dict[str, list]:
        return {"k": _K_GRID, **PARAM_GRIDS[self.basis]}
\end{lstlisting}

Consequências desse desenho:

\begin{itemize}
  \item Cada execução do pipeline já é, por construção, escopada a uma única base — não existe
        mais a possibilidade de misturar bases no mesmo resumo.
  \item \code{parameter\_grid()}/\code{random\_search\_space()} voltam a ser um único
        \code{Dict[str, list]} simples (\code{k} + os parâmetros \textbf{próprios} daquela base),
        sem precisar de um espaço de busca condicional/ramificado.
  \item \code{PARAM\_GRIDS}/\code{RANDOM\_SEARCH\_SPACES} — os dicionários que mapeiam nome de
        base para a faixa de busca de cada parâmetro — continuam centralizados aqui, e não dentro
        de \arq{core/math/basis/}, por decisão deliberada: são faixas de busca de
        \textbf{experimento} (decisão de quem desenha o experimento), não uma propriedade
        matemática da base (que pertence a \arq{core/}).
  \item \code{\_validate\_basis\_entry(basis\_name)} confere, na \textbf{construção} de cada
        \code{FEMaPlugin}, que \code{PARAM\_GRIDS}/\code{RANDOM\_SEARCH\_SPACES} cobrem
        \textbf{exatamente} os campos que \code{basis.PARAMS | basis.OPTIONAL\_PARAMS} declara —
        nem a mais, nem a menos. Isso significa que registrar uma base nova em \arq{core/} sem
        adicionar a entrada correspondente aqui produz um erro imediato e claro, no momento de
        instanciar o plugin, em vez de uma falha silenciosa ou um espaço de busca incompleto no
        meio de um tuning de horas.
\end{itemize}

O adaptador \code{\_FEMaEstimator} é o que dá ao FEMa uma interface ``\emph{sklearn-like}''
(\code{fit}/\code{predict}/\code{predict\_proba}) para caber no contrato \code{ModelPlugin}: ele
escolhe \code{FEMaClassifier} ou \code{FEMaRegressor} conforme o \code{context}, instancia a base
via \code{Basis.get}, e traduz \code{predict} de acordo (um classificador retorna
\code{(labels, probs)}; um regressor retorna só o vetor de previsões).

A função \code{available\_bases()} apenas repassa \code{Basis.available()} — usada por
\code{pipeline.run\_model.run\_all\_bases} e por \code{reporting/compare\_bases.py} para nunca
duplicar essa lista em outro lugar do projeto.

\section{\arq{logreg.py}, \arq{knn.py}, \arq{mlp.py} — baselines externos}

\code{models/registry.py} deixa claro, na própria docstring, que \textbf{o FEMa não está neste
registro}: ele não é identificado por um único nome, mas por um par \code{(context, basis)} (ver
\code{FEMaPlugin} acima), então não cabe no formato ``nome $\to$ instância fixa'' usado para
baselines. Este registro existe apenas para os baselines \textbf{externos}, que servem de
referência metodológica e não fazem parte da comparação de bases:

\begin{itemize}
  \item \code{LogRegPlugin} (\arq{logreg.py}) — regressão logística do scikit-learn, com cuidado
        explícito de só passar \code{penalty}/\code{solver} quando diferentes do padrão (evita um
        \code{FutureWarning} de depreciação em versões recentes do scikit-learn).
  \item \code{KNNPlugin} (\arq{knn.py}) — k-Vizinhos Mais Próximos
        (\code{KNeighborsClassifier}). É o baseline mais metodologicamente próximo do FEMa (também
        baseado em vizinhança), o que o torna uma referência particularmente informativa: mostra
        o quanto a \emph{escolha da base de interpolação} agrega sobre uma vizinhança
        ``crua''.
  \item \code{MLPPlugin} (\arq{mlp.py}) — uma rede neural simples
        (\code{MLPClassifier}) \textbf{existe como arquivo}, mas \textbf{não está registrada} em
        \code{models/registry.py::\_PLUGINS} na versão atual do código. Ou seja, o arquivo está
        pronto e funcional, mas inacessível via \code{model\_name="mlp"} até que alguém a adicione
        de volta ao dicionário \code{\_PLUGINS} — um único ponto de re-habilitação (ver
        Seção~\ref{sec:expansao-modelos}).
\end{itemize}

Ambos os baselines registrados (\code{logreg}, \code{knn}) seguem exatamente o mesmo padrão:
\code{parameter\_grid()} para grade discreta, \code{random\_search\_space()} para distribuições
contínuas, sem nenhuma lógica além da tradução direta para os construtores do scikit-learn.


\chapter{Métricas: \arq{metrics/}}

\section{\arq{contracts.py} — o contrato de métrica}

Toda métrica é uma função com a assinatura
\code{metric(y\_true, y\_pred, y\_score, n\_classes) -> Optional[float]}. \code{y\_score} é a
matriz de probabilidades (ou \code{None}, se indisponível); \code{n\_classes} é usado por métricas
que se comportam diferente em binário vs.\ multiclasse. Retornar \code{None} (em vez de levantar
exceção) é o sinal esperado para ``não calculável nesta situação'' — por exemplo, \code{roc\_auc}
sem \code{y\_score}, ou quando um fold contém uma única classe.

\section{\arq{classification.py} e \arq{regression.py} — implementações}

Sete métricas de classificação (via \code{sklearn.metrics}): \code{accuracy},
\code{balanced\_accuracy}, \code{precision}, \code{recall}, \code{f1}, \code{roc\_auc},
\code{mcc}. Todas são ``\emph{multiclass-aware}'': quando há mais de duas classes, usam
\code{average='macro'} (não ponderado por suporte, para não mascarar desempenho ruim em classes
minoritárias) e, no caso de \code{roc\_auc}, \code{multi\_class='ovr'}.

Cinco métricas de regressão: \code{mae}, \code{mse}, \code{rmse}, \code{r2}, \code{mape}. Seguem o
\textbf{mesmo contrato} (recebem \code{y\_score}/\code{n\_classes}, mas os ignoram) — isso é o que
permite que o mesmo motor de avaliação (\code{compute\_all}) sirva tanto para classificação quanto
para regressão sem nenhuma ramificação de código.

\section{\arq{registry.py} — registro central e direção de otimização}

\code{\_METRICS} mapeia nome $\to$ função; \code{HIGHER\_IS\_BETTER} mapeia nome $\to$ booleano
(usado por \code{persistence/summary\_builder.py} para saber se, ao rankear pela métrica
primária, um valor maior é melhor — caso de \code{f1} — ou pior — caso de \code{mae}). Duas listas
adicionais, \code{CLASSIFICATION\_METRICS} e \code{REGRESSION\_METRICS}, são a fonte única de
verdade para ``todas as métricas registradas'' de cada tipo de tarefa, acessadas através de
\code{metrics\_for\_context(context)}. Isso é usado por \code{pipeline/run\_model.py} para decidir
os \emph{defaults} de \code{metric\_names} conforme o contexto (classificador $\to$ 7 métricas;
regressor $\to$ 5 métricas), e por \code{reporting/compare\_bases.py} e pelos testes para garantir
que a comparação cobre \textbf{todas} as métricas relevantes, e não um subconjunto escolhido à mão
em cada lugar que precisa dessa lista.

\code{compute\_all(metric\_names, y\_true, y\_pred, y\_score, n\_classes)} calcula todas as
métricas pedidas de uma vez, capturando qualquer exceção individual e gravando \code{None} sem
interromper as demais.

\chapter{Busca de hiperparâmetros: \arq{tuning/}}

\section{\arq{grid\_search.py} — Grid Search}

\code{expand\_param\_grid(param\_grid)} expande um dicionário \code{\{chave: [valores]\}} no
produto cartesiano de todas as combinações. Também suporta, por compatibilidade histórica,
\code{List[Dict[str, list]]} — uma ``ramificação'' por dicionário, útil quando um hiperparâmetro só
faz sentido para certos valores de outro (o caso que motivou originalmente esse formato foi o
antigo \code{FEMaPlugin} genérico, descrito na Seção~\ref{sec:models-fema}; hoje nenhum plugin do
projeto usa mais esse formato condicional, mas o suporte permanece disponível para um plugin
futuro que precise dele). \code{combo\_id(params)} gera um identificador estável e legível (string
ordenada por chave), usado como chave de checkpoint e no ranking do resumo.

\section{\arq{random\_search.py} — Random Search}

\code{sample\_param\_space(search\_space, n\_iter, seed)} amostra \code{n\_iter} combinações de um
espaço de busca contínuo, de forma determinística dada a \code{seed}. Suporta os três tipos de
distribuição descritos na Seção~\ref{ch:models} (\code{uniform}, \code{loguniform},
\code{randint}) e também listas de escolha discreta. Assim como \code{grid\_search.py}, aceita o
formato condicional \code{List[Dict]} por compatibilidade — no caso condicional, cada amostra
sorteia primeiro \textbf{qual} ramificação usar (peso igual entre elas) e só então amostra os
campos daquela ramificação, nunca misturando campos de ramificações diferentes numa mesma amostra.

\section{\arq{param\_space.py} — o ponto único de entrada}

\code{generate\_combinations(plugin, strategy, seed, n\_iter, fixed\_hyperparameters)} é a única
função que \arq{pipeline/run\_model.py} chama: decide, com base em \code{strategy}
(\code{"grid\_search"} ou \code{"random\_search"}), qual dos dois motores usar. Se
\code{fixed\_hyperparameters} for informado, ignora completamente a estratégia de tuning e usa
exatamente aquela combinação — útil para reproduzir um resultado específico ou rodar sem busca de
hiperparâmetros (por exemplo, ao gerar uma curva de sensibilidade a um único parâmetro mantendo os
demais fixos).

\section{\arq{cv\_strategy.py} — geração de splits de validação cruzada}

\code{build\_cv\_splits(X, y, strategy, n\_splits, n\_repeats, master\_seed)} é um gerador que
produz tuplas \code{(repeat\_idx, fold\_idx, train\_idx, val\_idx)}. Duas estratégias:
\code{stratified\_kfold} (uma única passada) e \code{repeated\_stratified\_kfold} (repete com
seeds diferentes, reduzindo a variância da estimativa — útil com datasets pequenos ou
desbalanceados). \emph{Nested} CV não é oferecido de propósito: a separação treino/validação
(usada aqui para o tuning) e o \emph{holdout} de teste (Capítulo~\ref{ch:evaluation}) já cumprem o
mesmo papel metodológico sem o custo computacional multiplicado de uma validação cruzada
aninhada. A seed do particionador é sempre \textbf{derivada} da seed mestre
(\code{utils.seeding.derive\_seed}), garantindo que o mesmo conjunto de splits seja gerado sempre
que o pipeline for reexecutado com a mesma configuração — condição necessária para o checkpoint
funcionar corretamente (Capítulo~\ref{ch:persistence}).

\chapter{Treino: \arq{training/}}

Dois arquivos pequenos e sem lógica de negócio própria — só orquestram chamadas ao
\code{ModelPlugin}:

\begin{itemize}
  \item \code{trainer.py::train\_estimator(plugin, params, X\_train, y\_train, random\_state,
        X\_val, y\_val)} — instancia (\code{plugin.create\_model}) e treina
        (\code{plugin.fit}) um estimador. Não sabe nada sobre validação cruzada, busca de
        hiperparâmetros, métricas ou persistência.
  \item \code{final\_fit.py::fit\_final\_model(plugin, best\_params, X\_train, y\_train,
        master\_seed, X\_val, y\_val)} — chamado \textbf{depois} que o tuning + CV determina os
        melhores hiperparâmetros (Capítulo~\ref{ch:persistence}), treina um modelo
        \textbf{novo do zero} (não reaproveita nenhum modelo de fold), usando uma seed derivada
        especificamente para o retreino final (\code{derive\_seed(master\_seed, "final\_fit",
        plugin.name)}). Este é o único modelo que será avaliado no conjunto de teste.
\end{itemize}

\chapter{Avaliação: \arq{evaluation/}}
\label{ch:evaluation}

\begin{itemize}
  \item \code{cv\_evaluator.py::evaluate\_fold(plugin, params, X\_train\_fold, y\_train\_fold,
        X\_val\_fold, y\_val\_fold, metric\_names, random\_state)} — treina no split de treino do
        fold, avalia no split de validação, retorna as métricas calculadas. Não sabe nada sobre
        persistência.
  \item \code{test\_evaluator.py::evaluate\_on\_test(plugin, estimator, X\_test, y\_test,
        y\_train\_reference, metric\_names)} — chamado \textbf{uma única vez}, depois do
        retreino final. O conjunto de teste nunca é usado antes deste ponto — nem no
        grid/random search, nem na validação cruzada, nem em nenhuma decisão de tuning. Isso é o
        que garante que a métrica de teste reportada é uma estimativa honesta de generalização, e
        não um número inflado por vazamento de informação do teste para o tuning.
\end{itemize}

Ambos delegam o cálculo efetivo das métricas a \code{metrics.registry.compute\_all}
(Capítulo anterior) — a única diferença entre eles é \emph{quando} e \emph{sobre qual partição de
dados} são chamados.

\chapter{Persistência: \arq{persistence/}}
\label{ch:persistence}

Esta camada sofreu uma mudança arquitetural deliberada e importante: \textbf{deixou de saber o que
é ``modelo'' ou ``experimento''}. As três funções centrais de \arq{run\_writer.py}
(\code{next\_run\_path}, \code{load\_all\_runs}, \code{write\_summary}, \code{write\_test\_results})
recebem um número arbitrário de segmentos de caminho, \code{*scope}, e simplesmente os concatenam
sob \code{results/}:

\begin{lstlisting}
def _run_dir(*scope: str) -> Path:
    d = RESULTS_DIR.joinpath(*scope)
    d.mkdir(parents=True, exist_ok=True)
    return d
\end{lstlisting}

O \textbf{significado} de cada segmento — se é um contexto, uma base, um dataset, um nome de
experimento, ou o nome de um baseline — é responsabilidade exclusiva de quem monta o caminho
(\arq{pipeline/run\_model.py}, Capítulo~\ref{ch:pipeline}), nunca da camada de persistência. Essa
inversão de dependência é o que permite que \code{results/} seja organizado por
\textbf{base} (o eixo científico do projeto) sem que \arq{persistence/} precise ter qualquer
conhecimento sobre o que é uma ``base''.

\section{Escrita atômica}

\code{write\_run\_atomic(path, payload)} grava em um arquivo temporário no \textbf{mesmo}
diretório de destino e só então faz \code{os.replace} (rename atômico ao nível do sistema
operacional). Isso garante que nunca existirá, em disco, um JSON parcialmente escrito — mesmo que
o processo seja interrompido (\texttt{Ctrl+C}, queda de energia, \emph{out-of-memory}) no meio da
escrita.

\section{Checkpoint: retomada de execuções interrompidas}

\code{checkpoint.py::get\_completed\_keys(*scope)} varre todos os \code{run\_*.json} já existentes
em um diretório de escopo e reconstrói o conjunto de chaves \code{(combo\_id, repeat\_idx,
fold\_idx)} já executadas. \code{is\_done(completed\_keys, combo\_id, repeat\_idx, fold\_idx)} é a
checagem usada, fold a fold, dentro do laço principal do pipeline (Capítulo~\ref{ch:pipeline}),
para pular o que já foi feito. Isso é o que torna o pipeline \textbf{idempotente}: chamá-lo duas
vezes com os mesmos argumentos não duplica trabalho nem resultados — a segunda chamada
simplesmente constata que tudo já está pronto e vai direto para o resumo.

\section{Resumo: \arq{summary\_builder.py}}

\code{build\_summary(*scope, primary\_metric, higher\_is\_better, extra\_fields)} lê todos os
\code{run\_*.json} de um escopo, agrupa por \code{combo\_id}, calcula média e desvio-padrão de
cada métrica entre folds/repetições, e ordena as combinações pela métrica primária (respeitando a
direção de otimização — Capítulo anterior). O resultado (\code{summary.json}) inclui
\code{best\_configuration} (a melhor combinação encontrada) e o \code{ranking} completo.

Como consequência direta da mudança para \code{FEMaPlugin(context, basis)} fixo
(Seção~\ref{sec:models-fema}), \textbf{cada diretório de resultados corresponde hoje a uma única
base do FEMa} (ou a um único baseline externo) — logo, o ranking dentro de um \code{summary.json}
é sempre \textbf{homogêneo}: só varia \code{k} e os parâmetros próprios daquela base. A comparação
\emph{entre} bases diferentes é responsabilidade de \arq{reporting/} (Capítulo~\ref{ch:reporting}),
que lê um \code{summary.json} por base.

\code{extra\_fields} é injetado no resumo exatamente como veio de quem chamou (por exemplo,
\code{\{"context": "classifier", "basis": "shepard", "experiment": "baseline"\}}) —
\code{summary\_builder.py} não precisa interpretar o que esses campos significam, apenas
propagá-los para o JSON final.

\chapter{Orquestração: \arq{pipeline/run\_model.py}}
\label{ch:pipeline}

Este é o módulo que \textbf{liga todas as camadas anteriores}. Reflete, na própria API pública, a
divisão conceitual do projeto entre dois tipos de experimento com objetivos completamente
diferentes:

\begin{itemize}
  \item \textbf{\code{run\_basis\_experiment(context, basis, dataset, \ldots)}} — roda o FEMa com
        o contexto e a base fixos. É o pipeline \textbf{central} da pesquisa: cada chamada
        corresponde a uma célula da matriz (contexto $\times$ base) sendo avaliada. Persiste em
        \code{results/<context>/<basis>/<dataset>/<experiment\_name>/}.
  \item \textbf{\code{run\_all\_bases(context, dataset, bases=None, \ldots)}} — chama
        \code{run\_basis\_experiment} para \textbf{todas} as bases registradas em
        \code{core.Basis.available()} (ou para um subconjunto, se \code{bases} for informado). É o
        modo de uso mais comum do projeto: comparar todas as bases de uma vez, para um mesmo
        contexto/dataset.
  \item \textbf{\code{run\_baseline\_experiment(model\_name, dataset, \ldots)}} — roda um baseline
        externo (\code{logreg}, \code{knn}). \textbf{Não} faz parte da comparação de bases; serve
        apenas de referência metodológica (``o FEMa com a melhor base é competitivo com um
        classificador simples?''). Persiste, isolado, em
        \code{results/external\_baselines/<model\_name>/<dataset>/<experiment\_name>/}.
\end{itemize}

Os três compartilham o mesmo motor interno, \code{\_execute(scope, plugin, dataset, \ldots)}, que
implementa o fluxo completo:

\begin{enumerate}
  \item Fixa a seed global (\code{utils.seeding.set\_global\_seed}).
  \item Carrega o dataset (\code{datasets.loader.load\_dataset}).
  \item Gera as combinações de hiperparâmetros a avaliar (\code{tuning.param\_space.generate\_combinations}).
  \item Consulta o checkpoint (\code{persistence.checkpoint.get\_completed\_keys}) para saber o que
        já foi feito.
  \item Para cada combinação $\times$ cada fold de validação cruzada ainda não executado: treina e
        avalia (\code{evaluation.cv\_evaluator.evaluate\_fold}), mede o tempo de execução
        (\code{utils.timing.timer}), e grava o resultado imediatamente
        (\code{persistence.run\_writer.write\_run\_atomic}) — nunca espera o fim de tudo para
        persistir nada.
  \item Constrói o resumo (\code{persistence.summary\_builder.build\_summary}), que identifica a
        melhor combinação de hiperparâmetros.
  \item Retreina um modelo final com essa combinação, no conjunto de treino completo
        (\code{training.final\_fit.fit\_final\_model}).
  \item Avalia esse modelo final, uma única vez, no conjunto de teste
        (\code{evaluation.test\_evaluator.evaluate\_on\_test}), e grava o resultado
        (\code{persistence.run\_writer.write\_test\_results}).
\end{enumerate}

O motor \textbf{não sabe} se \code{plugin} é o FEMa ou um baseline externo — fala apenas com o
contrato \code{ModelPlugin} (Capítulo~\ref{ch:models}). A diferença entre
\code{run\_basis\_experiment} e \code{run\_baseline\_experiment} está inteiramente em \emph{como
cada uma monta o \code{scope} (caminho de resultados) e os \code{metric\_names}/\code{ranking\_metric}
padrão} antes de chamar \code{\_execute}.

Note, em particular, que \code{run\_basis\_experiment} escolhe as métricas e a métrica de
\emph{ranking} \textbf{padrão} de acordo com o contexto:

\begin{lstlisting}
is_classifier = context == "classifier"
default_metrics = DEFAULT_CLASSIFICATION_METRICS if is_classifier else DEFAULT_REGRESSION_METRICS
default_ranking = "f1" if is_classifier else "rmse"
\end{lstlisting}

Isso é o que torna a análise de regressão (Seção~\ref{sec:expansao-regressao}) uma questão de
\emph{configuração}, e não de código novo: bastam um dataset com alvo contínuo e
\code{context="regressor"}.

\section{\arq{run\_from\_experiment\_file} — roteamento de configuração}

L� \code{configs/base.yaml} (defaults globais) e o YAML do experimento (Capítulo~\ref{ch:configs}),
combina os dois, e decide qual das três funções acima chamar, com base no \emph{schema} do YAML:

\begin{itemize}
  \item Se o YAML contém a chave \code{basis} (uma base) ou \code{all\_bases: true} $\to$
        experimento do FEMa (\code{run\_basis\_experiment} ou \code{run\_all\_bases}). A chave
        \code{context} é opcional, com \emph{default} \code{"classifier"}.
  \item Se contém a chave \code{model} $\to$ baseline externo
        (\code{run\_baseline\_experiment}).
\end{itemize}

Isso significa que, na prática, o usuário quase nunca chama \code{run\_basis\_experiment} ou
\code{run\_baseline\_experiment} diretamente — escreve um arquivo YAML pequeno
(Capítulo~\ref{ch:configs}) e roda via \arq{main.py} (Capítulo~\ref{ch:main}).

\chapter{Relatórios e comparação de bases: \arq{reporting/}}
\label{ch:reporting}

\section{\arq{compare\_bases.py} — o módulo ativo}

Este é o módulo que materializa a pergunta de pesquisa do projeto: lê os \code{summary.json} e
\code{test\_results.json} de cada base já rodada (via
\code{pipeline.run\_model.run\_basis\_experiment}/\code{run\_all\_bases}) e produz, num único
ponto de entrada:

\begin{lstlisting}
from reporting.compare_bases import run_full_comparison

run_full_comparison(
    context="classifier", dataset="fetal_health", experiment_name="baseline",
    output_dir="reports/basis_comparison/classifier/fetal_health",
)
\end{lstlisting}

\begin{itemize}
  \item \code{comparison\_table.csv} — uma linha por base, uma coluna por métrica de teste
        (na ordem canônica de \code{metrics.registry.metrics\_for\_context}), mais a melhor
        combinação de hiperparâmetros de cada base.
  \item \code{comparison.json} — o mesmo conteúdo, em JSON, mais um \code{ranking\_per\_metric}:
        para cada métrica, as bases ordenadas da melhor para a pior (respeitando
        \code{HIGHER\_IS\_BETTER}).
  \item \code{plots/bar\_<metrica>.png} — um gráfico de barras por métrica, todas as bases lado a
        lado, gerado para \textbf{toda} métrica registrada para o contexto que tenha ao menos um
        valor disponível.
  \item \code{plots/curve\_<metrica>\_vs\_k.png} — uma curva de sensibilidade a $k$ por base
        (uma linha por base), construída a partir do \code{ranking} de validação cruzada de cada
        \code{summary.json} — \textbf{não} do teste, já que o teste só produz um único ponto por
        base, insuficiente para uma curva. A própria docstring da função documenta uma
        limitação conhecida: como o \emph{random search} varia o parâmetro próprio de cada base
        (\code{z}, \code{epsilon}, \ldots) simultaneamente a \code{k}, os pontos de uma mesma
        base podem não formar uma curva estritamente monotônica em \code{k} — o valor no eixo
        $Y$ para um dado \code{k} é de \textbf{uma} combinação específica sorteada com aquele
        \code{k}, não uma média marginal sobre o parâmetro próprio da base. Para uma curva
        estritamente marginal em \code{k}, é preciso rodar um experimento com
        \code{fixed\_hyperparameters} variando apenas \code{k}.
\end{itemize}

Os resultados de exemplo já presentes no repositório
(\arq{reports/basis\_comparison/classifier/fetal\_health/}) foram gerados exatamente por esse
fluxo, sobre o dataset \code{fetal\_health}.

\section{\arq{compare\_models.py} — arquivo legado, não utilizado}

Este arquivo é uma versão \textbf{anterior} de \code{compare\_bases.py} — de fato, sua própria
docstring já afirma ``\emph{esta é a substituição de reporting/compare\_models.py (removido)}'',
mas o arquivo antigo, na prática, não foi removido do repositório, apenas duplicado com o novo
nome. \code{reporting/\_\_init\_\_.py} importa exclusivamente de \code{compare\_bases}, então
\code{compare\_models.py} \textbf{não é executado por nenhum caminho de código ativo do
projeto} atualmente. Recomenda-se removê-lo numa próxima limpeza, para não confundir alguém que o
encontre e não perceba que está obsoleto (ver Capítulo~\ref{ch:legado}).

\section{\arq{plots.py} e \arq{curves.py} — utilitários de visualização gerais}

Dois módulos independentes de \code{compare\_bases.py}, com funções de plotagem mais genéricas
(matriz de confusão, curva ROC/PR a partir de \code{y\_score}, gráfico de valor previsto vs.\ real
para regressão) e o cálculo manual (sem dependência além de \emph{numpy}) de curvas ROC e
precision-recall em \code{curves.py}. Não são chamados automaticamente por nenhum pipeline —
existem como ferramentas prontas para uso pontual/exploratório (por exemplo, dentro de um notebook
de análise).

\chapter{Utilitários: \arq{utils/}}

\begin{itemize}
  \item \code{logging\_config.py::get\_logger(name)} — logger único por módulo, formato
        padronizado, grava simultaneamente em \emph{stdout} e em \code{logs/pipeline.log}. Evita
        handlers duplicados em reimportações.
  \item \code{seeding.py} — \code{set\_global\_seed(seed)} fixa as sementes de \code{random} e
        \emph{numpy}; \code{derive\_seed(master\_seed, *parts)} deriva, via \emph{hash} SHA-256
        estável, uma seed determinística a partir de identificadores arbitrários (nome do plugin,
        \code{combo\_id}, índice de fold, índice de repetição, \ldots). Isso garante que o
        \textbf{mesmo} fold/repetição sempre receba a \textbf{mesma} seed entre execuções
        diferentes — condição necessária tanto para reprodutibilidade quanto para o checkpoint.
  \item \code{timing.py::timer()} — \emph{context manager} simples para medir tempo decorrido
        (usado para registrar \code{execution\_time\_seconds} em cada \code{run\_*.json}).
  \item \code{hashing.py::stable\_hash} / \code{run\_identity} — geração de hashes SHA-256
        determinísticos a partir de dicionários (canonicalizando chaves ordenadas e arredondando
        \emph{floats} para evitar que \code{0.1} vs.\ \code{0.100000000000001} gerem hashes
        diferentes). \textbf{Não é mais usado} pelo sistema de checkpoint atual (que usa
        \code{combo\_id} + índices de fold/repetição diretamente, ver
        Capítulo~\ref{ch:persistence}) — é um remanescente de uma versão anterior do checkpoint,
        que identificava cada execução por um hash único. Mantido no repositório por ser
        potencialmente reaproveitável, mas hoje um código morto do ponto de vista do pipeline
        principal.
  \item \code{hardware\_info.py::get\_hardware\_info} / \code{get\_code\_version} — coleta
        informação de plataforma, contagem de CPUs, disponibilidade de GPU/\emph{CUDA} (via
        \emph{torch}, se instalado) e o \emph{commit} \emph{git} atual. Também \textbf{não é
        chamado} por nenhum ponto do pipeline atual — outro remanescente disponível para uso
        futuro (por exemplo, anexar essas informações a cada \code{run\_*.json}, para auditar em
        que máquina/versão do código cada resultado foi gerado).
\end{itemize}

\chapter{Configuração declarativa: \arq{configs/}}
\label{ch:configs}

\section{\arq{base.yaml} — defaults globais}

\begin{lstlisting}[style=yaml]
master_seed: 42
cross_validation:
  strategy: repeated_stratified_kfold
  n_splits: 5
  n_repeats: 3
tuning:
  strategy: random_search
  n_iter: 20
metrics: [accuracy, f1, balanced_accuracy, mcc]
ranking_metric: f1
higher_is_better: true
run_final_test: true
\end{lstlisting}

Todo campo aqui pode ser sobrescrito por um YAML de experimento individual.

\section{\arq{experiments/*.yaml} — um arquivo por experimento}

Quatro exemplos já presentes ilustram os dois \emph{schemas} de experimento reconhecidos por
\code{run\_from\_experiment\_file} (Capítulo~\ref{ch:pipeline}):

\begin{lstlisting}[style=yaml]
# fema_baseline.yaml — roda TODAS as bases para um contexto/dataset
context: classifier
all_bases: true
dataset: fetal_health
experiment_name: baseline
tuning_strategy: random_search
tuning_n_iter: 20
\end{lstlisting}

\begin{lstlisting}[style=yaml]
# fema_shepard_only.yaml — escopado a UMA UNICA base
context: classifier
basis: shepard
dataset: fetal_health
experiment_name: baseline
\end{lstlisting}

\begin{lstlisting}[style=yaml]
# knn.yaml — baseline externo, fora da comparacao de bases
model: knn
dataset: fetal_health
experiment_name: baseline
\end{lstlisting}

\code{smoke\_test.yaml} roda o baseline \code{logreg} sobre o dataset sintético, com validação
cruzada reduzida (3 \emph{folds}, 1 repetição) e apenas 3 iterações de \emph{random search} —
existe para validar rapidamente que o pipeline inteiro (do carregamento de dados à avaliação no
teste) funciona de ponta a ponta, sem esperar o tempo de um experimento real.

\chapter{Ponto de entrada: \arq{main.py}}
\label{ch:main}

A interface de linha de comando cobre quatro modos de uso, que podem ser combinados:

\begin{lstlisting}[style=yaml]
# Roda um ou mais experimentos a partir de arquivo(s) YAML
python main.py --experiment configs/experiments/fema_baseline.yaml

# Roda todas as bases diretamente por flags (sem YAML)
python main.py --context classifier --all-bases --dataset fetal_health

# Roda uma unica base
python main.py --context classifier --basis shepard --dataset fetal_health

# Aciona a comparacao entre bases ja rodadas (nao roda experimentos)
python main.py --compare --context classifier --dataset fetal_health

# Roda todas as bases E compara em seguida, num unico comando
python main.py --context classifier --all-bases --dataset fetal_health --compare

# Baseline externo
python main.py --baseline-model knn --dataset fetal_health
\end{lstlisting}

A flag \code{--compare} não roda nenhum experimento por si só — ela lê os resultados já
persistidos em \code{results/<context>/<basis>/<dataset>/<experiment\_name>/} e gera a tabela, o
JSON consolidado e os gráficos via \code{reporting.compare\_bases.run\_full\_comparison}
(Capítulo~\ref{ch:reporting}). Pode ser combinada com \code{--all-bases} para, num único comando,
rodar todas as bases e já produzir a comparação em seguida.

\chapter{Testes: \arq{tests/}}

\section{\arq{test\_pipeline.py} — testes de fumaça do pipeline}

Usa o dataset sintético (rápido, sem depender de arquivos externos) para validar, de ponta a
ponta: que o FEMa treina e prediz corretamente (\code{test\_fema\_classifier\_end\_to\_end}, com
um exemplo geometricamente óbvio — dois agrupamentos bem separados); que
\code{run\_basis\_experiment} persiste no caminho correto,
\code{results/<context>/<basis>/<dataset>/<experiment\_name>/}; que
\code{run\_baseline\_experiment} persiste isolado em
\code{results/external\_baselines/\ldots}; e que o checkpoint efetivamente evita duplicar
execuções (roda o mesmo experimento duas vezes e confere que o número de arquivos
\code{run\_*.json} não muda na segunda vez).

\section{\arq{test\_reporting.py} — testes do módulo de comparação}

Roda \code{run\_all\_bases} sobre um pequeno subconjunto de bases (\code{shepard}, \code{radial},
\code{rbf\_gaussian}) no dataset sintético, usando \textbf{todas} as métricas registradas para o
contexto (não um subconjunto reduzido), e confere que: cada linha da comparação traz todas as
métricas esperadas; o JSON consolidado e o CSV contêm exatamente essas métricas; e
\code{run\_full\_comparison} gera um gráfico de barras \textbf{e} uma curva de $k$ para
\textbf{cada} métrica registrada — nenhuma a mais, nenhuma a menos.

\section{\arq{tests/linear\_algebra/} — validação matemática das bases}

Dois arquivos com sobreposição significativa de conteúdo: \code{math\_validation.py} (mais antigo)
e \code{validation.py} (mais recente e mais completo — ver Capítulo~\ref{ch:legado}). Ambos
implementam um \emph{harness} que verifica propriedades matemáticas de cada base registrada em
\code{Basis.available()}, automaticamente (sem precisar listar as bases manualmente): sanidade
estrutural (não levantar exceção nem retornar \code{NaN}), partição da unidade
($\sum w_i = 1$), não-negatividade, igualdade de pesos para distâncias iguais, monotonicidade, e
uma métrica específica do domínio chamada \textbf{\emph{leakage}} (vazamento de interpolação).

A ideia de \emph{leakage} merece destaque por ser conceitualmente central à validação: o que
importa para a qualidade da interpolação não é o vetor de pesos em si comparado a um ``vetor
ideal'' $[1, 0, 0, \ldots]$, mas sim o quanto a predição final \textbf{pode} se afastar do valor
verdadeiro no pior caso. Com partição da unidade e o vizinho de distância zero no índice $j$:
$$\text{predição} - y_j = \sum_i w_i (y_i - y_j) \quad\Rightarrow\quad
|\text{predição} - y_j| \le (1 - w_j)\max_i|y_i - y_j|.$$
O termo $(1-w_j)$ — a soma dos pesos em todos os vizinhos \textbf{exceto} o de distância zero — é
o \emph{leakage}: simultaneamente o limite superior do erro de interpolação e um valor atingível.
A versão mais recente (\code{validation.py}) vai além da dedução algébrica e verifica isso
\textbf{empiricamente} (\code{check\_leakage\_matches\_worst\_case\_prediction}): monta um caso
sintético de pior caso ($y_j=0$, $y_i=1$ nos demais) e confere que a predição real bate com a
fórmula fechada do \emph{leakage}.

\code{validation.py} também inclui \code{SWEEP\_CONFIG}: para cada base, qual campo de
\code{BasisParameters} é o ``parâmetro de escala'' (o que controla a concentração da função perto
de $d=0$) e em que faixa varrê-lo, permitindo auditar como o \emph{leakage} muda conforme esse
parâmetro varia — inclusive documentando, na própria configuração, os casos em que uma base
\textbf{não tem} hoje um parâmetro de escala (\code{lorentzian}, \code{attention}), o que é
descrito como uma limitação real da fórmula atual, não uma omissão do harness.

\chapter{Código legado e pendências}
\label{ch:legado}

Para que este documento seja útil como mapa fiel do repositório (e não apenas da arquitetura
``ideal''), vale registrar explicitamente o que existe no código hoje mas não faz parte do
caminho ativo:

\begin{itemize}
  \item \textbf{\arq{src/}} — a árvore de módulos anterior à arquitetura atual
        (\code{config.py}, \code{cross\_validation.py}, \code{curves.py}, \code{datasets.py},
        \code{metrics.py}, \code{models.py}, \code{persistence.py}, \code{pipeline.py},
        \code{plots.py}, \code{tuning.py}, \code{utils.py}). Nenhum módulo fora de \arq{src/}
        importa nada dela. É segura para remoção; foi mantida até aqui apenas como referência
        histórica de como o projeto era organizado antes da divisão em camadas descrita neste
        documento.
  \item \textbf{\arq{reporting/compare\_models.py}} — versão anterior de
        \code{compare\_bases.py} (Seção~\ref{ch:reporting}), não referenciada por
        \code{reporting/\_\_init\_\_.py} nem por nenhum outro módulo. Candidato a remoção.
  \item \textbf{\arq{models/mlp.py}} — implementado, mas não registrado em
        \code{models/registry.py::\_PLUGINS}. Re-habilitá-lo é uma mudança de uma linha (ver
        Seção~\ref{sec:expansao-modelos}).
  \item \textbf{\arq{utils/hashing.py}} e \textbf{\arq{utils/hardware\_info.py}} — utilitários
        funcionais, mas não chamados por nenhum ponto do pipeline atual (o checkpoint usa
        \code{combo\_id} + índices de fold/repetição, não hashes; nenhum \code{run\_*.json} grava
        informação de hardware hoje). Ficam disponíveis para uso futuro.
  \item \textbf{\arq{tests/linear\_algebra/math\_validation.py}} — sobreposto por
        \code{validation.py} (mais completo, com verificação empírica de \emph{leakage} e
        \code{SWEEP\_CONFIG}). Mantido, mas duplica boa parte do conteúdo do arquivo mais novo.
  \item \textbf{\arq{teste.ipynb}} — notebook de uso exploratório, fora do escopo desta
        documentação (não faz parte do pipeline reprodutível).
\end{itemize}

Nenhum destes pontos é urgente — o pipeline funciona corretamente com eles presentes —, mas vale
uma limpeza numa próxima iteração, para reduzir a chance de alguém editar por engano um arquivo
que não é mais executado.

\chapter{Como expandir o projeto}
\label{ch:expansao}

Este capítulo é o mais orientado à ação: cobre, com passos concretos, as quatro direções de
expansão mais prováveis para este projeto. Em todos os casos, a resposta curta é a mesma —
\textbf{a arquitetura em camadas foi desenhada exatamente para que essas expansões não exijam
tocar em \arq{core/}, \arq{pipeline/}, \arq{persistence/}, \arq{training/} ou
\arq{evaluation/}}. Cada extensão se resume a adicionar um arquivo pequeno (ou uma entrada em um
dicionário) na camada correspondente.

\section{Adicionando um novo dataset}
\label{sec:expansao-dataset}

Passos:

\begin{enumerate}
  \item Coloque o CSV em \arq{data/raw/}.
  \item Adicione uma entrada em \code{datasets/registry.py::\_REGISTRY}:
\begin{lstlisting}
"meu_dataset": DatasetSpec(
    name="meu_dataset",
    csv_path=DATA_RAW / "meu_arquivo.csv",
    target_column="minha_coluna_alvo",
    drop_columns=("id",),        # opcional
    delimiter=",",                 # opcional, default ja e' ","
),
\end{lstlisting}
  \item Pronto — \code{datasets.loader.load\_dataset("meu\_dataset")} já funciona, com split
        treino/validação/teste determinístico e padronização automática. Nenhum outro arquivo
        precisa ser tocado: qualquer experimento (FEMa ou baseline) passa a poder usar
        \code{dataset: meu\_dataset} no YAML (Capítulo~\ref{ch:configs}).
\end{enumerate}

Se o alvo for \textbf{contínuo} (dataset de regressão), não é preciso nenhum passo adicional
aqui — \code{loader.py} já detecta automaticamente se a coluna-alvo é numérica ou categórica.
A diferenciação entre classificação e regressão acontece na escolha do \code{context} ao rodar o
experimento (próxima seção), não no cadastro do dataset.

\section{Adicionando uma nova base de interpolação ao FEMa}
\label{sec:expansao-base}

Este é o fluxo de expansão mais frequente do projeto, e o mais bem pavimentado — a própria
existência de \code{Basis.available()} como fonte única de verdade (Seção~\ref{sec:factory}) foi
desenhada para isso.

\begin{enumerate}
  \item Crie \arq{core/math/basis/minha\_base.py}:
\begin{lstlisting}
import numpy as np
from .base_basis import BaseBasis, DELTA
from .parameters import BasisParameters

class MinhaBase(BaseBasis):
    """phi(d) = ... (descreva a formula e cite a referencia, se houver)."""
    PARAMS = ("epsilon",)  # os campos de BasisParameters que a formula usa

    def evaluate(self, dists: np.ndarray, params: BasisParameters) -> np.ndarray:
        epsilon = self._require(params)["epsilon"]
        dists = np.where(dists == 0, DELTA, dists)  # se houver risco de divisao por zero
        return ...  # a formula em si
\end{lstlisting}
        Não implemente \code{compute\_weights} — a normalização já está pronta em
        \code{BaseBasis} (Seção~\ref{sec:factory-basis}).
  \item Registre em \code{core/math/basis/factory\_basis.py}:
\begin{lstlisting}
from .minha_base import MinhaBase
_CLASSES = {
    ...,
    "minha_base": MinhaBase,
}
\end{lstlisting}
  \item Adicione a entrada correspondente em \code{core/math/basis/\_\_init\_\_.py}
        (\code{import} + \code{\_\_all\_\_}), para consistência com as demais bases.
  \item Em \code{models/fema.py}, adicione a faixa de busca de hiperparâmetros em
        \textbf{ambos} \code{PARAM\_GRIDS} e \code{RANDOM\_SEARCH\_SPACES}:
\begin{lstlisting}
PARAM_GRIDS["minha_base"] = {"epsilon": [0.1, 1.0, 5.0, 20.0]}
RANDOM_SEARCH_SPACES["minha_base"] = {"epsilon": ("loguniform", 0.1, 50.0)}
\end{lstlisting}
        Se esquecer este passo, \code{\_validate\_basis\_entry} levanta um \code{ValueError}
        claro assim que alguém tentar instanciar \code{FEMaPlugin(context=\ldots,
        basis="minha\_base")} — o erro aparece na hora certa (construção do plugin), não no meio
        de um tuning de horas.
  \item (Recomendado) Adicione uma entrada em \code{SWEEP\_CONFIG}, em
        \code{tests/linear\_algebra/validation.py}, para que a nova base seja incluída
        automaticamente na varredura de validação matemática (partição da unidade,
        monotonicidade, \emph{leakage}, \ldots — Capítulo~\ref{ch:legado}).
\end{enumerate}

A partir daqui, a nova base já aparece automaticamente em \code{Basis.available()}, e portanto:
já pode ser usada em \code{basis: minha\_base} num YAML de experimento
(Seção~\ref{sec:models-fema}); já é incluída quando alguém roda \code{all\_bases: true}; e já
aparece automaticamente nas tabelas e gráficos de \code{reporting/compare\_bases.py}
(Capítulo~\ref{ch:reporting}) — nenhum desses três lugares precisa ser editado manualmente.

\section{Habilitando análise de regressão}
\label{sec:expansao-regressao}

Este é talvez o ponto mais importante deste capítulo: \textbf{a análise de regressão já está
implementada} — \code{FEMaRegressor} existe em \arq{core/models/fema\_regressor.py}
desde a base do projeto, \code{FEMaPlugin} já aceita \code{context="regressor"}
(Seção~\ref{sec:models-fema}), e \code{metrics/regression.py} já define as cinco métricas de
regressão (\code{mae}, \code{mse}, \code{rmse}, \code{r2}, \code{mape}). O que falta para rodar
um experimento de regressão de ponta a ponta é apenas:

\begin{enumerate}
  \item Um dataset com alvo \textbf{contínuo} registrado em \code{datasets/registry.py}
        (Seção~\ref{sec:expansao-dataset}) — nenhum dos quatro datasets atuais tem alvo contínuo,
        então este é o único passo genuinamente necessário.
  \item Rodar (ou escrever um YAML) com \code{context: regressor}:
\begin{lstlisting}
from pipeline.run_model import run_all_bases
run_all_bases(context="regressor", dataset="meu_dataset_continuo")
\end{lstlisting}
        ou, via YAML:
\begin{lstlisting}[style=yaml]
context: regressor
all_bases: true
dataset: meu_dataset_continuo
experiment_name: baseline
\end{lstlisting}
\end{enumerate}

A partir daí, todo o resto do pipeline já se adapta automaticamente: \code{run\_basis\_experiment}
já troca as métricas e a métrica de \emph{ranking} padrão para as de regressão
(\code{DEFAULT\_REGRESSION\_METRICS}, \code{ranking\_metric="rmse"} — Capítulo~\ref{ch:pipeline});
\code{\_FEMaEstimator} já instancia \code{FEMaRegressor} em vez de \code{FEMaClassifier}
(Seção~\ref{sec:models-fema}); \code{evaluate\_fold}/\code{evaluate\_on\_test} já chamam
\code{metrics.registry.compute\_all} com as métricas de regressão corretas; e os resultados vão
para \code{results/regressor/<basis>/<dataset>/<experiment\_name>/}, paralelo à árvore de
classificação, sem colidir com ela. \code{reporting.compare\_bases.run\_full\_comparison(context=
"regressor", \ldots)} já produz a tabela e os gráficos comparando as bases também neste contexto.

Em suma: nenhuma linha de \arq{core/}, \arq{pipeline/}, \arq{metrics/}, \arq{models/} ou
\arq{reporting/} precisa ser alterada para habilitar regressão — a única lacuna real é a ausência,
hoje, de um dataset de regressão cadastrado.

\section{Comparando com outros modelos (baselines)}
\label{sec:expansao-modelos}

Existem dois casos distintos aqui, que vale diferenciar:

\subsection{Re-habilitar um baseline já implementado (MLP)}

Como observado no Capítulo~\ref{ch:legado}, \code{models/mlp.py} já existe e é funcional — só não
está no registro. Basta:

\begin{lstlisting}
# models/registry.py
from models.knn import KNNPlugin
from models.logreg import LogRegPlugin
from models.mlp import MLPPlugin          # adicionar este import

_PLUGINS = {
    "logreg": LogRegPlugin(),
    "knn": KNNPlugin(),
    "mlp": MLPPlugin(),                    # adicionar esta linha
}
\end{lstlisting}

A partir daí, \code{model: mlp} já funciona em qualquer YAML de experimento, exatamente como
\code{logreg}/\code{knn}.

\subsection{Adicionar um baseline totalmente novo}

Passos, usando como exemplo um \code{RandomForestClassifier} do scikit-learn:

\begin{enumerate}
  \item Crie \arq{models/random\_forest.py}:
\begin{lstlisting}
from typing import Any, Dict
from sklearn.ensemble import RandomForestClassifier
from models.base import ModelPlugin

class RandomForestPlugin(ModelPlugin):
    name = "random_forest"
    supports_proba = True

    def create_model(self, params: Dict[str, Any], random_state: int):
        return RandomForestClassifier(
            n_estimators=params.get("n_estimators", 200),
            max_depth=params.get("max_depth", None),
            random_state=random_state,
        )

    def parameter_grid(self) -> Dict[str, list]:
        return {"n_estimators": [100, 200, 400], "max_depth": [None, 5, 10, 20]}

    def random_search_space(self) -> Dict[str, Any]:
        return {"n_estimators": ("randint", 50, 500), "max_depth": [None, 5, 10, 20, 40]}
\end{lstlisting}
  \item Registre em \code{models/registry.py::\_PLUGINS}, como acima.
  \item Crie um YAML em \arq{configs/experiments/} com \code{model: random\_forest}, ou chame
        \code{pipeline.run\_model.run\_baseline\_experiment(model\_name="random\_forest",
        dataset=\ldots)} diretamente.
\end{enumerate}

Nenhuma outra camada precisa mudar: \code{training/}, \code{evaluation/},
\code{tuning/} e \code{persistence/} já funcionam para qualquer \code{ModelPlugin}, seja ele o
FEMa ou um baseline externo. A única diferença estrutural entre um baseline e o FEMa é que
baselines usam um único nome fixo (\code{model\_name}) no registro, enquanto o FEMa é
parametrizado por \code{(context, basis)} — mas isso é transparente para quem só quer adicionar
\emph{mais um} baseline de comparação.

Para incluir o novo baseline na mesma tabela comparativa das bases do FEMa, basta chamar
\code{reporting.compare\_bases.compare} para o baseline separadamente (ele fica em
\code{results/external\_baselines/random\_forest/\ldots}) e concatenar as linhas manualmente, ou
escrever uma pequena função de comparação combinada — isso é o único ponto em que a separação
deliberada entre ``comparação de bases'' e ``baselines externos'' (Seção~\ref{ch:reporting}) exige
um passo manual extra, por design: o projeto não quer que um baseline externo apareça,
por acidente, misturado no ranking de bases do FEMa.

\section{Criando um novo experimento (YAML)}

Para qualquer combinação já suportada pelo código (dataset $\times$ contexto $\times$ base, ou
dataset $\times$ baseline), criar um novo experimento é só escrever um YAML novo em
\arq{configs/experiments/}, seguindo um dos dois \emph{schemas} da
Seção~\ref{ch:configs} — nunca é necessário editar \arq{pipeline/run\_model.py} para isso. Campos
omitidos herdam de \code{configs/base.yaml}.

\section{Adicionando uma nova métrica}

\begin{enumerate}
  \item Escreva a função em \code{metrics/classification.py} ou \code{metrics/regression.py},
        seguindo o contrato de \code{metrics/contracts.py}
        (\code{metric(y\_true, y\_pred, y\_score, n\_classes) -> Optional[float]}).
  \item Registre em \code{metrics/registry.py}: uma entrada em \code{\_METRICS}, uma em
        \code{HIGHER\_IS\_BETTER}, e inclua o nome em \code{CLASSIFICATION\_METRICS} ou
        \code{REGRESSION\_METRICS}, conforme o caso.
\end{enumerate}

A partir daí, a métrica passa a estar disponível para \code{metric\_names} em qualquer YAML de
experimento, para \code{ranking\_metric}, e é automaticamente incluída nas tabelas e nos gráficos
de \code{reporting/compare\_bases.py} — nenhum desses três consumidores precisa ser editado.

\section{Checklist resumo}

\begin{center}
\small
\begin{longtable}{@{}p{4.3cm}p{9.5cm}@{}}
\toprule
\textbf{Quero...} & \textbf{Arquivo(s) a tocar} \\
\midrule
\endhead
Novo dataset & \code{datasets/registry.py} (uma entrada) \\
Nova base do FEMa & novo arquivo em \code{core/math/basis/} + \code{factory\_basis.py} +
\code{core/math/basis/\_\_init\_\_.py} + \code{models/fema.py} (\code{PARAM\_GRIDS} e
\code{RANDOM\_SEARCH\_SPACES}) \\
Análise de regressão & apenas um dataset de alvo contínuo (Seção~\ref{sec:expansao-dataset}) +
\code{context: regressor} na chamada/YAML \\
Re-habilitar o MLP & \code{models/registry.py} (duas linhas) \\
Novo baseline & novo arquivo em \code{models/} + \code{models/registry.py} \\
Novo experimento & novo YAML em \code{configs/experiments/} \\
Nova métrica & função em \code{metrics/classification.py} ou \code{regression.py} +
\code{metrics/registry.py} \\
\bottomrule
\end{longtable}
\end{center}

\chapter{Conclusão}

A arquitetura descrita neste documento é o resultado de uma decisão de design consistente,
repetida em cada camada: \textbf{isolar o que muda (experimentos, datasets, bases, métricas) do
que não muda (a matemática em \arq{core/}, o motor de execução em
\arq{pipeline/run\_model.py})}. Isso é o que torna o projeto extensível na prática, e não apenas
em teoria — como ilustrado pelo Capítulo~\ref{ch:expansao}, cada uma das direções de expansão mais
prováveis (novo dataset, nova base, regressão, novo baseline) se resume a um arquivo pequeno ou a
uma entrada em um dicionário, nunca a uma reescrita de camadas já existentes.

\end{document}
